% Options for packages loaded elsewhere
\PassOptionsToPackage{unicode}{hyperref}
\PassOptionsToPackage{hyphens}{url}
\PassOptionsToPackage{dvipsnames,svgnames,x11names}{xcolor}
\documentclass[
  11pt,
  a4paper,
]{article}
\usepackage{xcolor}
\usepackage[margin=25mm]{geometry}
\usepackage{amsmath,amssymb}
\setcounter{secnumdepth}{-\maxdimen} % remove section numbering
\usepackage{iftex}
\ifPDFTeX
  \usepackage[T1]{fontenc}
  \usepackage[utf8]{inputenc}
  \usepackage{textcomp} % provide euro and other symbols
\else % if luatex or xetex
  \usepackage{unicode-math} % this also loads fontspec
  \defaultfontfeatures{Scale=MatchLowercase}
  \defaultfontfeatures[\rmfamily]{Ligatures=TeX,Scale=1}
\fi
\usepackage{lmodern}
\ifPDFTeX\else
  % xetex/luatex font selection
\fi
% Use upquote if available, for straight quotes in verbatim environments
\IfFileExists{upquote.sty}{\usepackage{upquote}}{}
\IfFileExists{microtype.sty}{% use microtype if available
  \usepackage[]{microtype}
  \UseMicrotypeSet[protrusion]{basicmath} % disable protrusion for tt fonts
}{}
\makeatletter
\@ifundefined{KOMAClassName}{% if non-KOMA class
  \IfFileExists{parskip.sty}{%
    \usepackage{parskip}
  }{% else
    \setlength{\parindent}{0pt}
    \setlength{\parskip}{6pt plus 2pt minus 1pt}}
}{% if KOMA class
  \KOMAoptions{parskip=half}}
\makeatother
\usepackage{longtable,booktabs,array}
\newcounter{none} % for unnumbered tables
\usepackage{calc} % for calculating minipage widths
% Correct order of tables after \paragraph or \subparagraph
\usepackage{etoolbox}
\makeatletter
\patchcmd\longtable{\par}{\if@noskipsec\mbox{}\fi\par}{}{}
\makeatother
% Allow footnotes in longtable head/foot
\IfFileExists{footnotehyper.sty}{\usepackage{footnotehyper}}{\usepackage{footnote}}
\makesavenoteenv{longtable}
\setlength{\emergencystretch}{3em} % prevent overfull lines
\providecommand{\tightlist}{%
  \setlength{\itemsep}{0pt}\setlength{\parskip}{0pt}}
\usepackage{microtype}
\usepackage{needspace}
\usepackage{fvextra}
\usepackage{adjustbox}
\AtBeginEnvironment{longtable}{\small\setlength{\tabcolsep}{3pt}}
\DefineVerbatimEnvironment{Highlighting}{Verbatim}{breaklines,commandchars=\\\{\}}
\setlength{\emergencystretch}{3em}
\widowpenalty=10000
\clubpenalty=10000
\allowdisplaybreaks
\usepackage{bookmark}
\IfFileExists{xurl.sty}{\usepackage{xurl}}{} % add URL line breaks if available
\urlstyle{same}
\hypersetup{
  pdftitle={One Fold{,} Unique Contact and the Complete Two-Piece Rate--Distortion Function of Rank-One Complex-Projective Born Prediction},
  pdfauthor={Lluis Eriksson (Independent researcher)},
  colorlinks=true,
  linkcolor={blue},
  filecolor={Maroon},
  citecolor={Blue},
  urlcolor={blue},
  pdfcreator={LaTeX via pandoc}}

\title{One Fold, Unique Contact and the Complete Two-Piece
Rate--Distortion Function of Rank-One Complex-Projective Born
Prediction}
\author{Lluis Eriksson (Independent researcher)}
\date{9 September 2026 - Internal revision v004}

\begin{document}
\maketitle

\subsection{Abstract}\label{abstract}

For a Haar-random pure state on \(\mathbb{CP}^{d-1}\) with
squared-Frobenius distortion on density-matrix reproductions, the exact
rate--distortion function \(R_d(D)\) is known
(ARR-2026-1D2QYXPCVY9H7ANB) as a one-dimensional envelope over a radius
\(b\in[0,\allowbreak{}1]\) and a multiplier \(\lambda\); first contact was proved
there, but uniqueness of the positive branch and absence of later branch
exchange were left open (its Remark 3.2). We close that remark for every
finite \(d\ge3\), in a note written to be read without the earlier
record (only its envelope theorem is imported, with the elementary facts
it rests on reproved here). The tool is the first-order equation
\(\kappa M'=(\kappa-d+1)M+(d-1)\) for \(M={}_1F_1(1;\allowbreak{}d;\allowbreak{}\kappa)\), which
reduces the fold numerator \(K_d'-\kappa K_d''\) to the sign of
\(\kappa-\Psi_d(M)\) with \(\Psi_d\) an explicit rational function; a
transversality argument in the \((M,\allowbreak{}\kappa)\) plane then shows the
stationary multiplier \(\lambda(\kappa)\) has exactly one fold. This
yields a unique coexistence multiplier \(\lambda_{c,\allowbreak{}d}\), no reentrance,
and \(R_d=\) a linear face glued at \(D_{c,\allowbreak{}d}\) to an explicit
complex-Bingham curve. A computer-assisted theorem gives the one-change
law for the fold-numerator coefficients; \(d=2\) recovers the
Dytso--Cardone sphere result.

\subsection{1. Introduction}\label{introduction}

Let \(X\) be Haar-distributed on \(\mathbb{CP}^{d-1}\),
\(\rho_X=|x\rangle\langle x|\), and let a reproduction be any density
matrix \(\hat\rho\) produced by an arbitrary stochastic map of \(X\);
the cost is \(\|\rho_X-\hat\rho\|_F^2\) and
\(R_d(D)=\inf\{I(X;\allowbreak{}\hat\rho):\mathbb{E}\|\rho_X-\hat\rho\|_F^2\le D\}\).
Through the Born rule this is the classical information needed to
predict the outcome statistics of Haar (or 2-design) measurements to
squared error \(D/(d(d+1))\).

\textbf{Antecedents.} The record ARR-2026-1D2QYXPCVY9H7ANB (\emph{Exact
Rate--Distortion Theory for Complex-Projective Born Prediction:
Finite-Dimensional Bingham Frontiers, Thermodynamic Coexistence, and
Worst-State Capacity}; below ``1D2Q'') proved, for
\(0\le D\le R_{0,\allowbreak{}d}^2=1-1/d\), the exact envelope
\(R_d(D)=\sup_{\lambda\ge0}\{\lambda(R_{0,\allowbreak{}d}^2-D)-\max_{0\le b\le1}G_{d,\allowbreak{}\lambda}(b)\}\)
with \(G_{d,\allowbreak{}\lambda}(b)=K_d(2\lambda b)-\lambda R_{0,\allowbreak{}d}^2b^2\),
\(K_d=\log{}_1F_1(1;\allowbreak{}d;\allowbreak{}\cdot)-\kappa/d\) (its Theorem 2.1, whose spectral
input is the centred power-sum extremum of Brazitikos--Pandis {[}BP,
Prop. 5.5{]}); every positive active radius is attained by a covariant
complex-Bingham channel. Its Theorem 3.1 proved a discontinuous onset:
for \(d\ge3\) there is \(0<\lambda_{c,\allowbreak{}d}<\lambda_{0,\allowbreak{}d}=d(d+1)/2\) at
which \(b=0\) and \emph{some} \(b_{c,\allowbreak{}d}>0\) are both global maximizers.
Its Remark 3.2 explicitly does \textbf{not} assert uniqueness of the
positive maximizer for larger \(\lambda\), nor exclude a later branch
exchange; consequently the finite-\(d\) frontier was known only as an
envelope, not as an explicit two-piece curve, and the asymptotics
\(\lambda_{c,\allowbreak{}d}=\alpha_*d-c_*\log d+O(1)\) (1D2Q Thm. 7.3) described a
first contact, not \emph{the} contact. Two sibling records did reach the
complete statement in other geometries: ARR-2026-61Y0FFA39M8KMBJ5
(\emph{Complete Rank-Two Born-Prediction Rate--Distortion on
\(\mathrm{Gr}_{\mathbb C}(2,\allowbreak{}4)\)}; ``61Y0'', Lemma 6.1, Thm. 6.2, Cor.
6.3) via an explicit integer coefficient sequence with one sign change,
and ARR-2026-7H9FAPTBZA897AMJ (\emph{Complete Rate--Distortion Phase
Diagram of Haar Oriented Two-Planes: One-Change Hypergeometric
Coefficients, Unique Coexistence, and No Reentrance}; ``7H9F'', Thm.
8.2, Thm. 9.1) for oriented two-planes in \(\mathbb R^n\), \(n\ge6\),
via a parity-truncated binomial window and a likelihood-ratio ordering.
The 7H9F normalizer \(L_q(\kappa)=\mathbb E\cosh(\kappa T)\),
\(T\sim\mathrm{Beta}(1,\allowbreak{}q)\), is the even part of the \(\mathbb{CP}^{q}\)
normalizer \(M_{q+1}(\kappa)=\mathbb E e^{\kappa T}\), but its window
argument does not transfer to the full (non-parity-truncated) window
that appears here. Externally: for \(d=2\) the problem is the real
two-sphere and \(R_2\) is the Dytso--Cardone curve {[}DC{]}; generalized
Curie--Weiss models can have several first-order transitions
(Eisele--Ellis {[}EE{]}), so uniqueness is not a generic fact; and the
Kummer-ratio literature (Karp--Sitnik {[}KS{]}, Kalmykov--Karp {[}KK{]},
Sitnik {[}S{]}) studies monotonicity and Turán-type inequalities for
ratios such as \(M_{d+1}/M_d\), which is exactly \(1-b_d(\kappa)\)
below, but contains no statement about the fold of
\(\kappa/b_d(\kappa)\).

\textbf{Contribution.} (i) A first-order linear ODE for
\(M_d={}_1F_1(1;\allowbreak{}d;\allowbreak{}\kappa)\) (Lemma 3.1) that collapses the three-term
fold numerator to a rational expression in \((M,\allowbreak{}\kappa)\) (Lemma 3.2).
(ii) A phase-plane transversality theorem (Theorem 4.1): for every
\(d\ge3\) the stationary multiplier has exactly one fold, with the
proved bounds \(2(d-2)(d+1)/d\le\kappa_{f,\allowbreak{}d}<2d\) and
\(M_d(\kappa_{f,\allowbreak{}d})\ge d(d-1)/2\); no coefficient sign pattern is
needed. (iii) The complete radial phase and the two-piece
rate--distortion function at every finite \(d\ge3\) (Theorem 5.1),
closing 1D2Q Remark 3.2; the single statement imported from 1D2Q
(Theorem A) is stated explicitly, and the elementary facts about
\(\mathrm{Beta}(1,\allowbreak{}d-1)\) that 1D2Q used are reproved in Lemma 2.1. (iv)
The one-change coefficient law in the style of 61Y0/7H9F (Theorem 6.1):
proved analytically for every \(d\) and every index except two boundary
indices for \(3\le d\le35\), which are settled by an exact integer
check; for \(d\ge36\) an explicit tail bound closes them. (v) The
\(d=2\) corollary (no fold, continuous onset, Dytso--Cardone). Sections
7--9 give numerical tables, open questions, and reproducibility data.

\textbf{What this note is, and what a reader gains without 1D2Q.} This
is a standalone note, not a continuation: Section 2 restates every
object it uses (\(M_d\), \(K_d\), \(G_{d,\allowbreak{}\lambda}\), the stationary
parametrization (2.3)), proves in Lemma 2.1 the three elementary facts
about \(T\sim\mathrm{Beta}(1,\allowbreak{}d-1)\) that the earlier record used, and
imports exactly one statement from 1D2Q, Theorem A (the envelope formula
and the attainability of active radii), with the equation numbers of its
proof. Lemma 3.1 itself is not new as a formula: it is a contiguous
relation of Kummer's function specialized to first parameter \(1\),
where \({}_1F_1(0;\allowbreak{}d;\allowbreak{}\kappa)=1\) closes the recurrence into an
inhomogeneous first-order equation (cf.~{[}DLMF, §13.3{]}); we include
its two-line series proof. What is new is what it buys. (1) The
stationary radius becomes the elementary ratio \(b_d=1-M_{d+1}/M_d\) and
the three-term fold numerator \(K_d'-\kappa K_d''\) becomes a rational
function of \((M_d,\allowbreak{}\kappa)\) (Lemma 3.2), so the fold question is a sign
problem for \(\kappa-\Psi_d(M_d(\kappa))\) about a classical special
function, with no rate--distortion context needed. (2) That sign problem
is settled globally by a transversality argument in the \((M,\allowbreak{}\kappa)\)
plane (Theorem 4.1): the curve \(\kappa=\Psi_d(M)\) can be crossed by
the graph of \(\kappa(M)\) only in one direction on \((1,\allowbreak{}M_1)\) and only
in the other on \((M_1,\allowbreak{}\infty)\), which forces exactly one crossing and
gives explicit bounds on where it lies. This answers completely the
one-parameter Curie--Weiss-type problem ``maximize
\(K_d(2\lambda b)-\lambda(1-1/d)b^2\) over \(b\in[0,\allowbreak{}1]\)'' for every
\(\lambda\) (Theorem 5.1(a)--(c)), a problem in which several
first-order transitions are not excluded a priori {[}EE{]}. (3) The
explicit two-piece rate--distortion function (Theorem 5.1(f)) is then
the corollary of (2) and Theorem A. A reader who knows 1D2Q additionally
obtains the closure of its Remark 3.2 and the identification of its
first-contact asymptotics (Theorem 7.3 there) as \emph{the} contact.

\subsection{2. Setting and notation}\label{setting-and-notation}

Throughout \(d\ge2\) is an integer, \(\kappa>0\), and
\(T\sim\mathrm{Beta}(1,\allowbreak{}d-1)\) (density \((d-1)(1-t)^{d-2}\) on
\([0,\allowbreak{}1]\)). From 1D2Q eqs. (2)--(6):

\[M_d(\kappa)={}_1F_1(1;d;\kappa)=\sum_{j\ge0}\frac{\kappa^j}{(d)_j}=\mathbb E\,e^{\kappa T},\qquad K_d(\kappa)=\log M_d(\kappa)-\frac{\kappa}{d},\qquad R_{0,d}^2=1-\frac1d, \tag{2.1}\]

\[G_{d,\lambda}(b)=K_d(2\lambda b)-\lambda R_{0,d}^2b^2\ (0\le b\le1),\qquad c_d(\lambda)=-\lambda R_{0,d}^2+\max_{0\le b\le1}G_{d,\lambda}(b). \tag{2.2}\]

\(K_d\) is the cumulant generating function of \(T-1/d\), so
\(K_d'(\kappa)=\mathbb E_\kappa T-1/d\) and
\(K_d''(\kappa)=\mathrm{Var}_\kappa T>0\) under the tilted law
\(\propto e^{\kappa t}\). Write \(m_d=M_d'/M_d=K_d'+1/d\). The
stationary equation \(G_{d,\allowbreak{}\lambda}'(b)=0\) for \(b>0\) reads, with
\(\kappa=2\lambda b\), \(K_d'(\kappa)=R_{0,\allowbreak{}d}^2b\); hence the positive
stationary radii are parametrized by \(\kappa\):

\[b_d(\kappa)=\frac{K_d'(\kappa)}{R_{0,d}^2}=\frac{d\,m_d(\kappa)-1}{d-1},\qquad \lambda_d(\kappa)=\frac{\kappa}{2b_d(\kappa)}, \tag{2.3}\]

(1D2Q eq. (110)). Define the fold numerator, the coexistence function
and the polynomial-series numerator

\[\begin{aligned}H_d&=K_d'-\kappa K_d'',\qquad F_d=2K_d-\kappa K_d',\\ N_d&:=M_d^2H_d=M_dM_d'-\kappa M_dM_d''+\kappa M_d'^2-\frac{M_d^2}{d}\\&=\sum_{n\ge0}c_{d,n}\kappa^n.\end{aligned}\tag{2.4}\]

Then \(F_d'=H_d\), \(\lambda_d'=H_d/(2R_{0,\allowbreak{}d}^2b_d^2)\), and at a
stationary radius \(G_{d,\allowbreak{}\lambda}(b_d(\kappa))=F_d(\kappa)/2\) and
\(\partial_b^2G_{d,\allowbreak{}\lambda}=-(2\lambda/b)H_d(\kappa)\) (direct
differentiation; the three identities are rechecked symbolically in
{\small\texttt{sym\_\allowbreak{}check.\allowbreak{}py}}). We drop the index \(d\) when harmless.

\textbf{Theorem A (imported; 1D2Q Theorem 2.1 and its proof, eqs.
(106)--(112)).} Let \(d\ge2\). (A1) For \(0\le D\le R_{0,\allowbreak{}d}^2\),
\(R_d(D)=\sup_{\lambda\ge0}\{\lambda(R_{0,\allowbreak{}d}^2-D)-\max_{0\le b\le1}G_{d,\allowbreak{}\lambda}(b)\}\).
(A2) If \(\lambda>0\) and \(b\in(0,\allowbreak{}1]\) is a global maximizer of
\(G_{d,\allowbreak{}\lambda}\) (``active''), then with \(\kappa=2\lambda b\) the
covariant channel
\(dP_{X|U=u}/d\mu_d=e^{\kappa|\langle x,\allowbreak{}u\rangle|^2}/M_d(\kappa)\),
\(U\) Haar, has posterior mean \((1-b)I/d+b\rho_u\) and attains the pair
\(D_b=R_{0,\allowbreak{}d}^2(1-b^2)\),
\(I_b=\kappa m_d(\kappa)-\log M_d(\kappa)=\lambda(R_{0,\allowbreak{}d}^2-D_b)-\max_vG_{d,\allowbreak{}\lambda}(v)\).
(A3) Mixing channels with probabilities \(t_i\) independent of \(X\)
achieves distortion \(\sum_i t_iD_i\) and information at most
\(\sum_i t_iI_i\). Retaining the independent flag as an auxiliary output
gives equality for the information of that enlarged output; discarding
the flag cannot increase information. Equality for the rate-distortion
frontier used below follows from the matching converse bound. For
\(d\ge3\) the spectral input rests on Brazitikos--Pandis {[}BP,
Proposition 5.5{]} through 1D2Q Proposition A.1; it bounds centered
power sums at fixed Euclidean norm. The interface and the
information-theoretic steps needed here are detailed below; the cited
extremal inequality itself remains an external theorem. The three
elementary facts that 1D2Q also used (its eqs. (10), (11), (15)) are
reproved here:

\emph{Dependency and proof interface for Theorem A.} For a centered real
spectrum \(a\) with \(\sum a_i^2=1\), maximizers of a power sum
\(p_j(a)=\sum a_i^j\) have at most three distinct coordinates: their
Lagrange equation \(j x^{j-1}=\alpha+\beta x\) has at most two real
roots for odd \(j\ge3\) by convexity, and at most three for even
\(j\ge4\) by its derivative. The cases \(j=1,\allowbreak{}2\) are fixed by the
constraints. The grouped-coordinate inequality {[}BP, Proposition 5.5{]}
therefore gives \(|p_j(a)|\le p_j(w)\) for odd \(j\) and
\(p_j(a)\le p_j(w)\) for even \(j\), where \(w\) is the centered
unit-norm spectrum with one positive spike. Newton's expansion
\[h_n(a)=\sum_{m_1+2m_2+\cdots=n}\prod_{j=1}^n\frac{p_j(a)^{m_j}}{m_j!\,j^{m_j}}\]
implies \(h_n(a)\le h_n(w)\): bound each summand by its absolute value,
then use the nonnegative power sums of \(w\). Haar squared coordinates
have the \(\mathrm{Dirichlet}(1,\allowbreak{}\ldots,\allowbreak{}1)\) law, so for a Hermitian
\(A\) with spectrum \(a\),
\[\mathbb E_\mu e^{t x^*Ax}=\sum_{n\ge0}\frac{h_n(a)t^n}{(d)_n}.\] The
series is absolutely convergent; scaling the fixed norm extends the
comparison to every traceless \(A\) and \(t\ge0\). Thus at fixed
\(\|A\|_F=bR_{0,\allowbreak{}d}\) the Laplace transform is bounded by that of
\(A=b(\rho_u-I/d)\).

For any density-matrix reproduction \(y\), set \(A=y-I/d\) and
\(S_x=\rho_x-I/d\). Then \(0\le b=\|A\|_F/R_{0,\allowbreak{}d}\le1\), and the
preceding bound gives
\[Z_\lambda(y):=\int e^{-\lambda\|S_x-A\|_F^2}\,d\mu(x)\le \exp\{-\lambda R_{0,d}^2+G_{d,\lambda}(b)\}\le e^{c_d(\lambda)}.\]
Nonnegativity of relative entropy of \(P_{X|y}\) against
\(e^{-\lambda\|S_x-A\|^2}\mu/Z_\lambda(y)\) yields
\(I(X;\allowbreak{}y)+\lambda\mathbb E\|\rho_X-y\|_F^2\ge-c_d(\lambda)\). This proves
the lower bound in (A1). The finite-information case suffices; infinite
information is automatic.

Every positive active radius at finite \(\lambda\) is interior by
\(G'_{d,\allowbreak{}\lambda}(1)<0\) (H2), and hence \(K_d'(\kappa)=R_{0,\allowbreak{}d}^2b\). The
displayed Bingham posterior has mean \((1-b)I/d+b\rho_u\) by unitary
symmetry and \(\mathbb E_\kappa |\langle X,\allowbreak{}u\rangle|^2=m_d(\kappa)\).
The posterior-mean identity gives \(D_b=R_{0,\allowbreak{}d}^2(1-b^2)\), and direct
integration of its log density gives \(I_b=\kappa m_d-\log M_d\). The
map \(u\mapsto(1-b)I/d+b\rho_u\) is injective on projective space for
\(b>0\), so the same information is obtained with density-matrix output.
This proves (A2); independence of the mixture flag and data processing
give (A3). The radial analysis of Sections 4 and 5(a)-(e) uses no
envelope equality. Its active channels, together with (A3), attain the
lower bound for all distortions in Theorem 5.1(f) and Corollary 5.2,
completing (A1) without circularity. This verifies the use of the
earlier envelope theorem here while retaining {[}BP{]} as the stated
external mathematical input.

\emph{Clarification of the external power-sum input.} We use the
non-strict bound of {[}BP, Proposition 5.5{]}. Its arXiv v1 proof, p.31,
contains the erroneous estimate \(75/64<1\) for \(n=5\). The local
argument is repaired by retaining the exponent \(k\ge3\): with its
notation \(w^2=1/2\), \(s^2=4/5\) and \(q=t/w=1/\sqrt{10}<1/3\), the
finite geometric sum satisfies
\[\frac{S(w)}{s^{2k}}<\frac{2}{1-q}\left(\frac58\right)^k<3\left(\frac58\right)^3=\frac{375}{512}<1.\]
This restores that estimate for \(n=5,\allowbreak{}k\ge3\) and leaves the separately
treated quartic cases unchanged. For \(n=3,\allowbreak{}k=2\), the normalized fourth
moment is identically \(1/2\); the non-strict inequality required here
is unaffected. We import no uniqueness assertion for that exceptional
case. This local repair does not replace the remainder of the cited
proof.

\textbf{Lemma 2.1 (elementary facts on \(\mathrm{Beta}(1,\allowbreak{}d-1)\)).} Let
\(d\ge2\). (i) \(M_d(\kappa)=(d-1)\int_0^1e^{\kappa t}(1-t)^{d-2}\,\allowbreak{}dt\)
equals the series in (2.1), and \(M_d\) is entire. (ii)
\(\mathbb ET=1/d\), \(v_d:=\mathrm{Var}\,\allowbreak{}T=\dfrac{d-1}{d^2(d+1)}\),
\(\mu_{3,\allowbreak{}d}:=\mathbb E(T-1/d)^3=\dfrac{2(d-1)(d-2)}{d^3(d+1)(d+2)}\);
hence
\(K_d(\kappa)=\frac{v_d}{2}\kappa^2+\frac{\mu_{3,\allowbreak{}d}}{6}\kappa^3+O(\kappa^4)\)
and, with \(\kappa=2\lambda b\),
\[G_{d,\lambda}(b)=\lambda\big(2\lambda v_d-R_{0,d}^2\big)b^2+\tfrac43\lambda^3\mu_{3,d}\,b^3+O(b^4). \tag{2.5}\]
(iii)
\(\log M_d(\kappa)=\kappa-(d-1)\log\kappa+\log\Gamma(d)+\log P_d(\kappa)\)
with \(P_d(\kappa)=1-e^{-\kappa}\sum_{j=0}^{d-2}\kappa^j/j!\in(0,\allowbreak{}1)\)
and \(P_d(\kappa)\to1\); in particular
\(\log M_d(\kappa)=\kappa-(d-1)\log\kappa+\log\Gamma(d)+o(1)\) as
\(\kappa\to\infty\).

\emph{Proof.} (i) Termwise,
\((d-1)\int_0^1t^j(1-t)^{d-2}dt=(d-1)B(j+1,\allowbreak{}d-1)=j!\,\allowbreak{}(d-1)!/(d-1+j)!=j!/(d)_j\).
(ii) The same integral gives \(\mathbb ET^k=k!/(d)_k\), so
\(\mathbb ET=1/d\), \(\mathbb ET^2=2/(d(d+1))\),
\(\mathbb ET^3=6/(d(d+1)(d+2))\); expanding \(\mathbb E(T-1/d)^2\) and
\(\mathbb E(T-1/d)^3\) gives \(v_d\) and \(\mu_{3,\allowbreak{}d}\), and the second
and third cumulants of \(T-1/d\) are its second and third central
moments; (2.5) follows from
\(G_{d,\allowbreak{}\lambda}(b)=K_d(2\lambda b)-\lambda R_{0,\allowbreak{}d}^2b^2\). (iii)
Substituting \(t=1-s\) and \(u=\kappa s\),
\(M_d(\kappa)=(d-1)e^\kappa\kappa^{1-d}\int_0^\kappa e^{-u}u^{d-2}du=\Gamma(d)e^\kappa\kappa^{1-d}P_d(\kappa)\),
where \(\int_0^\kappa e^{-u}u^{d-2}du=(d-2)!\,\allowbreak{}P_d(\kappa)\) by repeated
integration by parts. \(\square\)

(These are 1D2Q eqs. (10), (11) and (15); the symbolic expansion (2.5)
is rechecked in {\small\texttt{h6\_\allowbreak{}remark63\_\allowbreak{}check.\allowbreak{}py}}.)

\subsection{3. Exact reductions}\label{exact-reductions}

\textbf{Lemma 3.1 (first-order equation).}
\(\kappa M_d'=(\kappa-d+1)M_d+(d-1)\). Consequently
\[b_d(\kappa)=1-\frac{d\,(M_d-1)}{\kappa M_d}=1-\frac{M_{d+1}(\kappa)}{M_d(\kappa)},\qquad \kappa K_d'(\kappa)=\kappa R_{0,d}^2-(d-1)+\frac{d-1}{M_d(\kappa)}. \tag{3.1}\]

\emph{Proof.} With \(a_j=1/(d)_j\) one has \(a_{j-1}=(d-1+j)a_j\), so
the \(\kappa^j\) coefficient of the right side is
\(a_{j-1}-(d-1)a_j=j\,\allowbreak{}a_j\) for \(j\ge1\) and \(-(d-1)+(d-1)=0\) for
\(j=0\), matching \(\kappa M'\). Dividing by \(\kappa M\) gives
\(m_d=1-(d-1)(M-1)/(\kappa M)\), and (2.3) gives the first form of
\(b_d\); the second follows from
\(d(M_d-1)/\kappa=\sum_{i\ge0}d\kappa^i/(d)_{i+1}=M_{d+1}\). Finally
\(\kappa K'=\kappa m_d-\kappa/d\). \(\square\)

Differentiating Lemma 3.1 gives Kummer's equation
\(\kappa M''+(d-\kappa)M'-M=0\), so no independent input is used below.

\textbf{Lemma 3.2 (rational reduction).} Put
\[\begin{aligned}\tilde N_d(\kappa)&:=(\kappa-2d)M_d^2+d(\kappa-d+3)M_d+d(d-1)\\&=M_d(M_d+d)\big(\kappa-\Psi_d(M_d)\big),\\ \Psi_d(M)&:=\frac{d(M-1)(2M+d-1)}{M(M+d)}.\end{aligned}\tag{3.2}\]
Then \(N_d=\dfrac{d-1}{d\,\allowbreak{}\kappa}\tilde N_d\); hence
\(\operatorname{sign}H_d(\kappa)=\operatorname{sign}\big(\kappa-\Psi_d(M_d(\kappa))\big)\).

\emph{Proof.} Let \(P=\kappa M'=AM+(d-1)\) with \(A=\kappa-d+1\);
Kummer's equation gives \(\kappa^2M''=\kappa M-(d-\kappa)P\). Then
\(\kappa N_d=MP-M(\kappa M-(d-\kappa)P)+P^2-\kappa M^2/d=(2-A)MP+P^2-\kappa(1+1/d)M^2\).
Expanding \(P\):
\((2-A)M(AM+d-1)+(AM+d-1)^2=2AM^2+(d-1)(A+2)M+(d-1)^2\), and
\(2A-\kappa(d+1)/d=\frac{d-1}{d}(\kappa-2d)\),
\((d-1)(A+2)=\frac{d-1}{d}\,\allowbreak{}d(\kappa-d+3)\),
\((d-1)^2=\frac{d-1}{d}\,\allowbreak{}d(d-1)\). The factorization in (3.2) is an
identity of rational functions (checked in {\small\texttt{sym\_\allowbreak{}check.\allowbreak{}py}}), and
\(M_d>1\) on \(\kappa>0\). \(\square\)

\emph{Remark 3.3.} Substituting \(\kappa=\) from (3.1) one also finds
\(\tilde N_d=\kappa M_d\,\allowbreak{}[(2b_d-1)M_d+(d-1)b_d+1]\), so \(H_d>0\)
wherever \(b_d\ge1/2\); every zero of \(H_d\) has \(b_d<1/2\).

\textbf{Proposition 3.4 (closed coefficients).} For \(n\ge0\) let
\(T=2d-2+n\), \(u_n=\binom{T-1}{d-2}\),
\(W_{n-1}=\sum_{k=d-1}^{d-2+n}\binom{T-1}{k}\) (so \(W_{-1}=0\)). Then
\[[\kappa^n]\tilde N_d=\frac{(d-1)!^2}{T!}\,\hat c_{d,n},\qquad \hat c_{d,n}=(n-2d-2)\,W_{n-1}+\frac{d}{d-1}\,n(n+2d)\,u_n, \tag{3.3}\]
and \(c_{d,\allowbreak{}n}=\frac{d-1}{d}\frac{(d-1)!^2}{(2d-1+n)!}\hat c_{d,\allowbreak{}n+1}\).
Moreover \(\hat c_{d,\allowbreak{}0}=\hat c_{d,\allowbreak{}1}=\hat c_{d,\allowbreak{}2}=0\) and
\(\hat c_{d,\allowbreak{}3}=-\frac{(2d+1)(d-2)}{(d+1)(d+2)}\binom{2d}{d}\),
i.e.~\([\kappa^3]\tilde N_d=-\frac{d-2}{d^2(d+1)(d+2)}\).

\emph{Proof.} With \(a_j=(d-1)!/(d-1+j)!\),
\([M^2]_n=\sum_{i+j=n}a_ia_j=\frac{(d-1)!^2}{T!}\sum_{k=d-1}^{d-1+n}\binom{T}{k}=\frac{(d-1)!^2}{T!}W_n\).
For \(n\ge1\),
\([\kappa^n]\tilde N_d=[M^2]_{n-1}-2d[M^2]_n+d\,\allowbreak{}a_{n-1}+d(3-d)a_n\);
multiplying by \(T!/(d-1)!^2\) gives
\(TW_{n-1}-2dW_n+dT\binom{T-1}{d-1}+d(3-d)\binom{T}{d-1}\). Pascal's
rule and the symmetry \(\binom{T-1}{d-1+n}=\binom{T-1}{d-2}\) give
\(W_n=2W_{n-1}+2u_n\), while \(\binom{T-1}{d-1}=u_n\frac{d-1+n}{d-1}\)
and \(\binom{T}{d-1}=u_n\frac{T}{d-1}\). Collecting, the \(W_{n-1}\)
coefficient is \(T-4d=n-2d-2\) and the \(u_n\) coefficient is
\(\frac{d}{d-1}[T(n+2)-4(d-1)]=\frac{d}{d-1}n(n+2d)\). The case \(n=0\)
is \(-2d+d(3-d)+d(d-1)=0\). The special values follow from
\(W_0=\binom{2d-2}{d-1}\), \(W_1=2\binom{2d-1}{d-1}\),
\(W_2=\binom{2d}{d}\frac{3d+1}{d+1}\) and
\(u_3=\binom{2d}{d}\frac{d(d-1)}{(d+1)(d+2)}\). \(\square\)

Formula (3.3) is verified against the three-term definition in exact
rational arithmetic for \(2\le d\le40\), \(0\le n\le400\)
({\small\texttt{verify\_\allowbreak{}signs.\allowbreak{}py}}) and for \(2\le d\le12\), \(n\le60\)
({\small\texttt{check\_\allowbreak{}formula.\allowbreak{}py}}).

\subsection{\texorpdfstring{4. One fold for every
\(d\ge3\)}{4. One fold for every d\textbackslash ge3}}\label{one-fold-for-every-dge3}

\textbf{Theorem 4.1.} Let \(d\ge3\). \(H_d\) has exactly one positive
zero \(\kappa_{f,\allowbreak{}d}\); \(H_d<0\) on \((0,\allowbreak{}\kappa_{f,\allowbreak{}d})\) and \(H_d>0\)
on \((\kappa_{f,\allowbreak{}d},\allowbreak{}\infty)\). Moreover
\[M_d(\kappa_{f,d})\ \ge\ M_1:=\frac{d(d-1)}2,\qquad \frac{2(d-2)(d+1)}{d}=\Psi_d(M_1)\ \le\ \kappa_{f,d}\ <\ 2d. \tag{4.1}\]
Hence \(\lambda_d\) decreases strictly on \((0,\allowbreak{}\kappa_{f,\allowbreak{}d}]\) from
\(\lambda_{0,\allowbreak{}d}=d(d+1)/2\) (its limit at \(0^+\)) to
\(\lambda_{\min,\allowbreak{}d}:=\lambda_d(\kappa_{f,\allowbreak{}d})\) and increases strictly on
\([\kappa_{f,\allowbreak{}d},\allowbreak{}\infty)\) to \(+\infty\).

\emph{Proof.} \(\kappa\mapsto M_d(\kappa)\) is analytic and strictly
increasing (\(M'=\mathbb E[Te^{\kappa T}]>0\)), a bijection
\((0,\allowbreak{}\infty)\to(1,\allowbreak{}\infty)\); let \(\kappa(M)\) be its \(C^\infty\)
inverse and \(D(M):=\kappa(M)-\Psi_d(M)\). By Lemma 3.2,
\(\operatorname{sign}H_d(\kappa)=\operatorname{sign}D(M_d(\kappa))\), so
it suffices to show that \(D\) has exactly one zero \(M_f\) on
\((1,\allowbreak{}\infty)\), with \(D<0\) before and \(D>0\) after. By Lemma 3.1,
\(d\kappa/dM=1/M'(\kappa)=\Phi(M,\allowbreak{}\kappa):=\kappa/(\kappa M-(d-1)(M-1))\),
so \[D'(M)=\Phi(M,\kappa(M))-\Psi_d'(M).\] At a zero \(M_0\) of \(D\),
\[D'(M_0)=\Delta(M_0):=\Phi(M_0,\Psi_d(M_0))-\Psi_d'(M_0).\] The
following are rational identities in \((M,\allowbreak{}d)\), obtained by direct
computation ({\small\texttt{sym\_\allowbreak{}check.\allowbreak{}py}}):

\begin{enumerate}
\def\labelenumi{(\alph{enumi})}
\item
  On the curve \(\kappa=\Psi_d(M)\):
  \(\kappa M-(d-1)(M-1)=\frac{(d+1)M(M-1)}{M+d}>0\) for \(M>1\), and
  \(\Phi(M,\allowbreak{}\Psi_d(M))=\frac{d(2M+d-1)}{(d+1)M^2}\).
\item
  \(\Psi_d'(M)=\frac{d[(d+3)M^2+2(d-1)M+d(d-1)]}{M^2(M+d)^2}>0\) for
  \(M\ge1\).
\item
  \(\Delta(M)=\dfrac{d\,\allowbreak{}(M-1)^2\,\allowbreak{}(2M-d(d-1))}{(d+1)M^2(M+d)^2}\):
  negative on \((1,\allowbreak{}M_1)\), zero at \(M_1=d(d-1)/2\), positive on
  \((M_1,\allowbreak{}\infty)\).
\item
  \(\partial_\kappa\Phi(M,\allowbreak{}\kappa)=-\frac{(d-1)(M-1)}{(\kappa M-(d-1)(M-1))^2}<0\)
  for \(M>1\) wherever the denominator is nonzero.
\item
  \(D<0\) on \((1,\allowbreak{}1+\varepsilon)\) for some \(\varepsilon>0\): by
  Proposition 3.4,
  \(\tilde N_d(\kappa)=-\frac{d-2}{d^2(d+1)(d+2)}\kappa^3+O(\kappa^4)<0\)
  for small \(\kappa>0\) when \(d\ge3\),
  i.e.~\(\kappa<\Psi_d(M_d(\kappa))\).
\item
  \(2d-\Psi_d(M)=\frac{d[(d+3)M+d-1]}{M(M+d)}>0\), so \(\Psi_d<2d\) and
  \(D(M)\to+\infty\) as \(M\to\infty\).
\end{enumerate}

\emph{Step 1: \(D<0\) on \((1,\allowbreak{}M_1)\) and \(D(M_1)\le0\).} Suppose
\(D(M)\ge0\) for some \(M\in(1,\allowbreak{}M_1)\) and let
\(M_*=\inf\{M\in(1,\allowbreak{}M_1):D(M)\ge0\}\). By (e), \(M_*\ge1+\varepsilon\);
by continuity \(D(M_*)=0\) and \(D<0\) on \((1,\allowbreak{}M_*)\), so
\(D'(M_*)\ge0\). But \(D'(M_*)=\Delta(M_*)<0\) by (c). Contradiction;
continuity gives \(D(M_1)\le0\).

\emph{Step 2: on \([M_1,\allowbreak{}\infty)\), \(D(M)<0\) implies \(D'(M)>0\).} If
\(D(M)<0\) then \(\kappa(M)<\Psi_d(M)\). The denominator of
\(\Phi(M,\allowbreak{}\cdot)\) is affine in \(\kappa\), equals
\(\kappa M'(\kappa)>0\) at \(\kappa=\kappa(M)\) and is positive at
\(\kappa=\Psi_d(M)\) by (a); hence it is positive on the whole segment
\([\kappa(M),\allowbreak{}\Psi_d(M)]\), where (d) applies, so
\(\Phi(M,\allowbreak{}\kappa(M))>\Phi(M,\allowbreak{}\Psi_d(M))\) and \(D'(M)>\Delta(M)\ge0\) by
(c).

\emph{Step 3: exactly one zero.} Let \(M_f=\inf\{M\ge M_1:D(M)\ge0\}\),
which exists by (f); then \(D(M_f)=0\) and, by Step 1 and the
definition, \(D<0\) on \((1,\allowbreak{}M_f)\). We claim \(D>0\) on
\((M_f,\allowbreak{}\infty)\). If \(D(M_2)<0\) for some \(M_2>M_f\), let
\(M_3=\sup\{M\in[M_f,\allowbreak{}M_2]:D(M)\ge0\}\); then \(D(M_3)=0\) and \(D<0\) on
\((M_3,\allowbreak{}M_2]\), so by Step 2 \(D\) is strictly increasing there and
\(D(M_2)>D(M_3)=0\), a contradiction. Thus \(D\ge0\) on
\([M_f,\allowbreak{}\infty)\). If \(D(M_0)=0\) for some \(M_0>M_f\ge M_1\), then
\(M_0\) is a local minimum of \(D\), so \(D'(M_0)=0\), whereas
\(D'(M_0)=\Delta(M_0)>0\) by (c). Hence \(D>0\) on \((M_f,\allowbreak{}\infty)\).
(When \(D(M_1)=0\) this gives \(M_f=M_1\); the argument does not need to
distinguish the cases.)

Therefore \(\kappa_{f,\allowbreak{}d}:=\kappa(M_f)\) is the unique positive zero of
\(H_d\), with the stated signs. Since \(M_f\ge M_1\) and
\(\kappa_{f,\allowbreak{}d}=\Psi_d(M_f)\) with \(\Psi_d\) increasing (b) and bounded
by \(2d\) (f), (4.1) follows; \(\Psi_d(M_1)=2(d-2)(d+1)/d\) is a direct
evaluation. The monotonicity of \(\lambda_d\) follows from
\(\lambda_d'=H_d/(2R_{0,\allowbreak{}d}^2b_d^2)\),
\(\lambda_d(0^+)=R_{0,\allowbreak{}d}^2/(2K_d''(0))=R_{0,\allowbreak{}d}^2/(2v_d)=d(d+1)/2\), and
\(\lambda_d(\kappa)\ge\kappa/2\to\infty\). \(\square\)

\textbf{Corollary 4.2 (\(d=2\)).} \(H_2>0\) on \((0,\allowbreak{}\infty)\);
\(\lambda_2\) increases strictly from \(\lambda_{0,\allowbreak{}2}=3\) to
\(+\infty\).

\emph{Proof.} Now \(M_1=1\), so \(\Delta>0\) on \((1,\allowbreak{}\infty)\) by (c).
Proposition 3.4 gives \(\hat c_{2,\allowbreak{}m}=(m-6)(2^{m+1}-2)+2m(m+4)\), so
\(\hat c_{2,\allowbreak{}3}=0\), \(\hat c_{2,\allowbreak{}4}=4\), and
\(\tilde N_2(\kappa)=\kappa^4/180+O(\kappa^5)>0\): \(D>0\) near \(1^+\).
If \(D\) had a zero, the first one, \(M_*\), would satisfy
\(D'(M_*)\le0\) (approached from \(D>0\)) and \(D'(M_*)=\Delta(M_*)>0\).
\(\square\)

\subsection{5. The chain to the rate--distortion
function}\label{the-chain-to-the-ratedistortion-function}

Standing facts, valid for \(d\ge2\) and all \(\kappa>0\), \(\lambda>0\):

(H1) \(b_d'=K_d''/R_{0,\allowbreak{}d}^2>0\); \(b_d(0^+)=0\); \(b_d(\kappa)\to1\) as
\(\kappa\to\infty\) (by (3.1), since \(d(M-1)/(\kappa M)\to0\)).

(H2) \(K_d'(\kappa)=\mathbb E_\kappa T-1/d<R_{0,\allowbreak{}d}^2\) because \(T<1\)
a.s.; hence \(G_{d,\allowbreak{}\lambda}'(1)=2\lambda(K_d'(2\lambda)-R_{0,\allowbreak{}d}^2)<0\),
so \(b=1\) is never stationary and never a maximizer.

(H3) For \(b\in(0,\allowbreak{}1]\): \(G_{d,\allowbreak{}\lambda}'(b)=0\iff b=b_d(\kappa)\) and
\(\lambda=\lambda_d(\kappa)\) with \(\kappa=2\lambda b\); the map
\(\kappa\mapsto b_d(\kappa)\) is injective by (H1).

(H4) At such a stationary radius \(G_{d,\allowbreak{}\lambda}(b)=F_d(\kappa)/2\) and
\(G_{d,\allowbreak{}\lambda}''(b)=-(2\lambda/b)H_d(\kappa)\).

(H5) \(F_d'=H_d\), \(F_d(0)=0\), and by (3.1) and Lemma 2.1(iii),
\(F_d(\kappa)=R_{0,\allowbreak{}d}^2\kappa-2(d-1)\log\kappa+O(1)\to+\infty\).

(H6) \(G_{d,\allowbreak{}\lambda}(0)=0\), \(G'_{d,\allowbreak{}\lambda}(0)=0\),
\(G''_{d,\allowbreak{}\lambda}(0)=2\lambda(2\lambda v_d-R_{0,\allowbreak{}d}^2)\) and
\(G'''_{d,\allowbreak{}\lambda}(0)=8\lambda^3\mu_{3,\allowbreak{}d}\), by (2.5). Since
\(R_{0,\allowbreak{}d}^2/(2v_d)=d(d+1)/2=\lambda_{0,\allowbreak{}d}\): for
\(\lambda<\lambda_{0,\allowbreak{}d}\) the radius \(0\) is a strict local maximum of
\(G_{d,\allowbreak{}\lambda}\), and for \(\lambda>\lambda_{0,\allowbreak{}d}\) a strict local
minimum. At \(\lambda=\lambda_{0,\allowbreak{}d}\) the quadratic term vanishes and
the dimensions differ: for \(d\ge3\), \(\mu_{3,\allowbreak{}d}>0\), so
\(G_{d,\allowbreak{}\lambda_{0,\allowbreak{}d}}(b)=\frac43\lambda_{0,\allowbreak{}d}^3\mu_{3,\allowbreak{}d}b^3+O(b^4)>0\)
for all small \(b>0\) and \(0\) is \emph{not} a local maximum; for
\(d=2\), \(\mu_{3,\allowbreak{}2}=0\) and
\(K_2(\kappa)=\log\big(\sinh(\kappa/2)/(\kappa/2)\big)\) gives
\(G_{2,\allowbreak{}3}(b)=-\frac{9}{20}b^4+\frac{9}{35}b^6+O(b^8)\), so \(0\) remains
a strict local maximum at \(\lambda=\lambda_{0,\allowbreak{}2}=3\) (in agreement with
Corollary 5.2, where \(0\) is the unique global maximizer for every
\(\lambda\le3\)). Only the cases \(\lambda<\lambda_{0,\allowbreak{}d}\),
\(\lambda>\lambda_{0,\allowbreak{}d}\) and (\(d\ge3\), \(\lambda=\lambda_{0,\allowbreak{}d}\)) are
used below.

\textbf{Theorem 5.1 (complete radial phase and two-piece RDF,
\(d\ge3\)).} Let \(\kappa_f=\kappa_{f,\allowbreak{}d}\),
\(\lambda_{\min}=\lambda_d(\kappa_f)\) be as in Theorem 4.1.

\begin{enumerate}
\def\labelenumi{(\alph{enumi})}
\item
  \emph{Stationary radii.} For \(\lambda>0\) the positive stationary
  radii of \(G_{d,\allowbreak{}\lambda}\) are the \(b_d(\kappa)\) with
  \(\lambda_d(\kappa)=\lambda\): none for \(\lambda<\lambda_{\min}\);
  exactly one, \(b_d(\kappa_f)\), for \(\lambda=\lambda_{\min}\); two,
  \(b_d(\kappa_-)<b_d(\kappa_+)\) with \(\kappa_-<\kappa_f<\kappa_+\),
  for \(\lambda_{\min}<\lambda<\lambda_{0,\allowbreak{}d}\); exactly one,
  \(b_d(\kappa_+)\) with \(\kappa_+>\kappa_f\), for
  \(\lambda\ge\lambda_{0,\allowbreak{}d}\). A pre-fold radius (\(\kappa<\kappa_f\))
  is a strict local minimum, a post-fold radius (\(\kappa>\kappa_f\)) a
  strict local maximum, and the degenerate radius at
  \(\lambda=\lambda_{\min}\) (where \(H_d(\kappa_f)=0\)) has gain
  \(G=F_d(\kappa_f)/2<0=G_{d,\allowbreak{}\lambda}(0)\).
\item
  \emph{Unique contact.} \(F_d\) decreases strictly on \((0,\allowbreak{}\kappa_f]\),
  is negative there, then increases strictly to \(+\infty\); it has a
  unique positive zero \(\kappa_c=\kappa_{c,\allowbreak{}d}>\kappa_f\). Put
  \(b_c=b_d(\kappa_c)\), \(\lambda_c=\lambda_d(\kappa_c)\),
  \(D_c=R_{0,\allowbreak{}d}^2(1-b_c^2)\),
  \(R_c=\kappa_cK_d'(\kappa_c)-K_d(\kappa_c)\).
\item
  \emph{Global maximizers.} The set of global maximizers of
  \(G_{d,\allowbreak{}\lambda}\) on \([0,\allowbreak{}1]\) is \(\{0\}\) for
  \(0<\lambda<\lambda_c\); \(\{0,\allowbreak{}b_c\}\) for \(\lambda=\lambda_c\);
  \(\{b_d(\kappa_+(\lambda))\}\) for \(\lambda>\lambda_c\), where
  \(\kappa_+(\lambda)>\kappa_c\) is the unique post-fold solution of
  \(\lambda_d(\kappa)=\lambda\). (At \(\lambda=0\), \(G_{d,\allowbreak{}0}\equiv0\)
  and every \(b\in[0,\allowbreak{}1]\) is a maximizer.)
\item
  \(0<\lambda_{\min}<\lambda_c<\lambda_{0,\allowbreak{}d}\).
\item
  \emph{No reentrance.} On \((\lambda_c,\allowbreak{}\infty)\),
  \(\lambda\mapsto b_d(\kappa_+(\lambda))\) is continuous and strictly
  increasing to \(1\), so the active distortion
  \(D(\lambda)=R_{0,\allowbreak{}d}^2(1-b_d(\kappa_+(\lambda))^2)\) decreases
  strictly from \(D_c\) to \(0\). The function \(c_d\) of (2.2) is
  convex with \(-\partial c_d(\lambda)=\{R_{0,\allowbreak{}d}^2\}\) for
  \(0<\lambda<\lambda_c\), \([D_c,\allowbreak{}R_{0,\allowbreak{}d}^2]\) at \(\lambda_c\),
  \(\{D(\lambda)\}\) for \(\lambda>\lambda_c\).
\item
  \emph{Two-piece rate--distortion function.}
  \[R_d(D)=\begin{cases}+\infty,& D=0,\\[2pt] \kappa K_d'(\kappa)-K_d(\kappa), & 0<D<D_c,\\[2pt] \lambda_c\,(R_{0,d}^2-D), & D_c\le D\le R_{0,d}^2,\\[2pt] 0,& D\ge R_{0,d}^2,\end{cases}\]
  where, on \(0<D<D_c\), \(\kappa>\kappa_c\) is the unique solution of
  \(b_d(\kappa)=\sqrt{1-D/R_{0,\allowbreak{}d}^2}\) and
  \(\kappa K_d^{\prime}(\kappa)-K_d(\kappa)=\kappa m_d(\kappa)-\log M_d(\kappa)\).
  The middle piece is attained by the covariant complex-Bingham channel
  of field \(\kappa\) and the linear face by an independently flagged
  mixture of the trivial channel and the field-\(\kappa_c\) channel. The
  pieces meet at \(D_c\) with value
  \(R_c=\lambda_cR_{0,\allowbreak{}d}^2b_c^2=\kappa_cR_{0,\allowbreak{}d}^2b_c/2\).
\end{enumerate}

\emph{Proof.} (a) By (H3) the positive stationary radii at \(\lambda\)
correspond bijectively to solutions \(\kappa>0\) of
\(\lambda_d(\kappa)=\lambda\), and Theorem 4.1 gives the counts: on
\((0,\allowbreak{}\kappa_f]\), \(\lambda_d\) is strictly decreasing with values in
\([\lambda_{\min},\allowbreak{}\lambda_{0,\allowbreak{}d})\) (the value \(\lambda_{0,\allowbreak{}d}\) is a
limit, not attained); on \([\kappa_f,\allowbreak{}\infty)\) it is strictly increasing
onto \([\lambda_{\min},\allowbreak{}\infty)\). The local type follows from (H4) and
the sign of \(H_d\); at \(\kappa_f\), (H4) and (b) give
\(G=F_d(\kappa_f)/2<0\).

\begin{enumerate}
\def\labelenumi{(\alph{enumi})}
\setcounter{enumi}{1}
\item
  \(F_d'=H_d\) has the sign pattern of Theorem 4.1, \(F_d(0)=0\), and
  (H5). So \(F_d<0\) on \((0,\allowbreak{}\kappa_f]\) and \(F_d\) increases strictly
  from \(F_d(\kappa_f)<0\) to \(+\infty\) on \([\kappa_f,\allowbreak{}\infty)\): one
  zero \(\kappa_c>\kappa_f\), with \(F_d<0\) on \((\kappa_f,\allowbreak{}\kappa_c)\)
  and \(F_d>0\) on \((\kappa_c,\allowbreak{}\infty)\).
\item
  Fix \(\lambda>0\). \(G_{d,\allowbreak{}\lambda}\) is continuous on \([0,\allowbreak{}1]\), so a
  maximizer exists among
  \(\{0\}\cup\{1\}\cup\{\text{interior stationary points}\}\); \(b=1\)
  is excluded by (H2), pre-fold radii by (a) (a strict local minimum of
  a nonconstant analytic function is not a global maximum), and the
  degenerate radius at \(\lambda_{\min}\) by (a). So the only
  competitors of \(0\) (gain \(0\)) are post-fold radii, present iff
  \(\lambda>\lambda_{\min}\), unique, equal to
  \(b_d(\kappa_+(\lambda))\) with gain \(F_d(\kappa_+(\lambda))/2\) by
  (H4). Since \(\kappa_+\) is the inverse of the strictly increasing
  \(\lambda_d|_{(\kappa_f,\allowbreak{}\infty)}\), it is continuous and strictly
  increasing; by (b) the gain is \(<0\), \(=0\), \(>0\) according as
  \(\kappa_+(\lambda)<\kappa_c\), \(=\kappa_c\), \(>\kappa_c\),
  i.e.~\(\lambda<\lambda_c\), \(=\lambda_c\), \(>\lambda_c\) (note
  \(\lambda_c=\lambda_d(\kappa_c)>\lambda_d(\kappa_f)=\lambda_{\min}\),
  so the range \(\lambda\le\lambda_{\min}\) is included in
  \(\lambda<\lambda_c\)).
\item
  \(\lambda_{\min}<\lambda_c\) was just shown. By (H6), at
  \(\lambda=\lambda_{0,\allowbreak{}d}\) small positive radii have positive gain, so
  \(0\) is not a global maximizer there; by (c) this forces
  \(\lambda_{0,\allowbreak{}d}>\lambda_c\).
\item
  The first sentence follows from (c), (H1) and the monotonicity of
  \(\kappa_+\). For each fixed \(b\),
  \(\lambda\mapsto K_d(2\lambda b)-\lambda R_{0,\allowbreak{}d}^2(1+b^2)\) is convex
  (\(K_d\) convex), so \(c_d\) is convex as a maximum of convex
  functions, and Danskin's theorem gives
  \(\partial c_d(\lambda)=\operatorname{conv}\{2bK_d'(2\lambda b)-R_{0,\allowbreak{}d}^2(1+b^2): b\text{ active}\}\);
  at an active \(b>0\), \(K_d'(\kappa)=R_{0,\allowbreak{}d}^2b\) turns the derivative
  into \(-R_{0,\allowbreak{}d}^2(1-b^2)=-D_b\), and at \(b=0\) it is \(-R_{0,\allowbreak{}d}^2\)
  (1D2Q (112)). Insert the active sets from (c).
\item
  \(R_d(D)=0\) for \(D\ge R_{0,\allowbreak{}d}^2\) (trivial channel) and
  \(R_d(0)=+\infty\) (zero distortion forces \(\hat\rho=\rho_X\) a.s.,
  and \(I(X;\allowbreak{}X)=+\infty\) for the continuous \(X\); quantitatively, 1D2Q
  Prop. 3.3). Let \(0<D<D_c\). By (H1) and (c)--(e), \(b_d\) maps
  \((\kappa_c,\allowbreak{}\infty)\) increasingly onto \((b_c,\allowbreak{}1)\), so there is a
  unique \(\kappa>\kappa_c\) with \(b_d(\kappa)=\sqrt{1-D/R_{0,\allowbreak{}d}^2}\);
  put \(\lambda=\lambda_d(\kappa)>\lambda_c\). By (c), \(b_d(\kappa)\)
  is the unique global maximizer of \(G_{d,\allowbreak{}\lambda}\), hence active, and
  \(\kappa=2\lambda b_d(\kappa)\). By (A2) the covariant channel attains
  distortion \(D\) at information
  \(I=\kappa m_d(\kappa)-\log M_d(\kappa)=\lambda(R_{0,\allowbreak{}d}^2-D)-\max_vG_{d,\allowbreak{}\lambda}(v)\).
  The right side is the expression inside the supremum of (A1) at this
  \(\lambda\), so \(R_d(D)\ge I\); attainment gives \(R_d(D)\le I\).
  Finally
  \(\kappa m_d-\log M_d=\kappa(K_d'+1/d)-K_d-\kappa/d=\kappa K_d'-K_d\).
  Let \(D_c\le D\le R_{0,\allowbreak{}d}^2\). Taking \(\lambda=\lambda_c\) in (A1),
  where \(\max G_{d,\allowbreak{}\lambda_c}=0\) by (c), gives
  \(R_d(D)\ge\lambda_c(R_{0,\allowbreak{}d}^2-D)\). Conversely both \(0\) and \(b_c\)
  are active at \(\lambda_c\); by (A2) the field-\(\kappa_c\) channel
  attains \((D_c,\allowbreak{}R_c)\) with \(R_c=\lambda_c(R_{0,\allowbreak{}d}^2-D_c)\) and the
  trivial channel attains \((R_{0,\allowbreak{}d}^2,\allowbreak{}0)\); by (A3) their flagged
  mixture with weight \(t=(R_{0,\allowbreak{}d}^2-D)/(R_{0,\allowbreak{}d}^2-D_c)\) achieves
  distortion \(D\) with information at most \(\lambda_c(R_{0,\allowbreak{}d}^2-D)\);
  the preceding lower bound forces equality. The two expressions for
  \(R_c\) agree because \(F_d(\kappa_c)=0\) means
  \(K_d(\kappa_c)=\kappa_cK_d'(\kappa_c)/2\), and
  \(\lambda_cR_{0,\allowbreak{}d}^2b_c^2=\kappa_cR_{0,\allowbreak{}d}^2b_c/2=\kappa_cK_d'(\kappa_c)/2\).
  \(\square\)
\end{enumerate}

\emph{Consequences.} The multiplier \(\lambda_{c,\allowbreak{}d}\) of 1D2Q Theorem
3.1 is the \(\lambda_c\) above, its coexisting positive radius is unique
and equals \(b_c\), and the asymptotics of 1D2Q Theorem 7.3 describe
\emph{the} contact. This closes 1D2Q Remark 3.2 for every finite
\(d\ge3\). By (4.1), \(\kappa_f=2d-O(1)\) and, by Remark 3.3,
\(b_d(\kappa_f)<1/2\), so
\(\lambda_{\min}=\kappa_f/(2b_d(\kappa_f))>\kappa_f\); numerically
\(\lambda_{\min}-2d\to0\), matching the birth of stationary points at
\(\alpha=2\) in 1D2Q Theorem 5.1.

\textbf{Corollary 5.2 (\(d=2\)).} For \(0<\lambda\le3\) the unique
global maximizer of \(G_{2,\allowbreak{}\lambda}\) is \(0\); for \(\lambda>3\) it is
the unique positive stationary radius, which tends to \(0\) as
\(\lambda\downarrow3\) (continuous onset). Consequently
\(R_2(D)=\kappa K_2'(\kappa)-K_2(\kappa)\) with
\(b_2(\kappa)=\sqrt{1-2D}\) for \(0<D<1/2\), which is the Dytso--Cardone
formula \(R_2=\eta r-\log(\sinh\eta/\eta)\), \(r=\coth\eta-1/\eta\),
under \(\kappa=2\eta\).

\emph{Proof.} By Corollary 4.2, \(\lambda_2(\kappa)>3\) for all
\(\kappa>0\) and \(\lambda_2\) is a bijection
\((0,\allowbreak{}\infty)\to(3,\allowbreak{}\infty)\). For \(\lambda\le3\), \(G'_{2,\allowbreak{}\lambda}\) has
no zero in \((0,\allowbreak{}1]\) and \(G'_{2,\allowbreak{}\lambda}(1)<0\) (H2), so
\(G_{2,\allowbreak{}\lambda}\) is strictly decreasing and \(0\) is the unique
maximizer. For \(\lambda>3\) there is exactly one stationary
\(b_+\in(0,\allowbreak{}1)\); \(G''(0)>0\) by (H6) and \(G'(1)<0\) give \(G'>0\) on
\((0,\allowbreak{}b_+)\) and \(G'<0\) on \((b_+,\allowbreak{}1]\). Theorem A applies verbatim (it
holds for \(d=2\)), and the proof of (f) goes through with \(D_c\)
replaced by \(R_{0,\allowbreak{}2}^2=1/2\). Finally
\(M_2(\kappa)=(e^\kappa-1)/\kappa\) gives
\(K_2(\kappa)=\log(\sinh\eta/\eta)\), \(b_2=2K_2'=\coth\eta-1/\eta\) and
\(\kappa K_2'=\eta r\). \(\square\)

\subsection{6. The coefficient sign law}\label{the-coefficient-sign-law}

Theorem 4.1 did not use the coefficients \(c_{d,\allowbreak{}n}\). The following
statement, in the format of 61Y0 Lemma 6.1 and 7H9F Theorem 8.2, is
nevertheless of independent interest and gives a second (Descartes-type)
proof of one fold.

\textbf{Theorem 6.1 (one-change law).} Let \(d\ge3\) and
\(N_d=\sum_nc_{d,\allowbreak{}n}\kappa^n\) as in (2.4). Then \(c_{d,\allowbreak{}0}=c_{d,\allowbreak{}1}=0\),
\[c_{d,n}<0\ \ (2\le n\le n_d),\qquad c_{d,n}>0\ \ (n>n_d),\qquad n_d=\begin{cases}2d-1,&d=3,4,\\ 2d,&d\ge5.\end{cases}\]
For \(d=2\): \(c_{2,\allowbreak{}0}=c_{2,\allowbreak{}1}=c_{2,\allowbreak{}2}=0\) and \(c_{2,\allowbreak{}n}>0\) for
\(n\ge3\). \emph{Status:} proved for all \(d\ge2\) and all \(n\), except
that for \(3\le d\le35\) the two indices \(n\in\{2d-1,\allowbreak{}2d\}\) are settled
by an exact integer computation ({\small\texttt{tail\_\allowbreak{}cases.\allowbreak{}py}}, which covers
\(3\le d\le60\)); for \(d\ge36\) these indices are covered by the
analytic bound in Lemma 6.2(iv).

By Proposition 3.4,
\(\operatorname{sign}c_{d,\allowbreak{}n}=\operatorname{sign}\hat c_{d,\allowbreak{}n+1}\); we
work with \(m=n+1\). For \(3\le m\le2d+1\) put
\[q_m:=\frac{W_{m-1}}{u_m},\qquad g_m:=\frac{d\,m(m+2d)}{(d-1)(2d+2-m)}.\]
Since \(m-2d-2<0\) for \(m\le2d+1\), (3.3) gives
\(\hat c_{d,\allowbreak{}m}<0\iff q_m>g_m\) for \(m\le2d+1\), while for \(m\ge2d+2\)
both terms of (3.3) are nonnegative and the second is positive, so
\(\hat c_{d,\allowbreak{}m}>0\).

\textbf{Lemma 6.2.} Let \(d\ge3\).

\begin{enumerate}
\def\labelenumi{(\roman{enumi})}
\item
  \emph{Recursion.} \(q_{m+1}=\dfrac{2(d+m)}{2d-2+m}(q_m+1)\) for
  \(m\ge1\), and \(q_3=\dfrac{(d+2)(3d+1)}{d(d-1)}\).
\item
  \emph{Base.} \(q_3-g_3=\dfrac{(2d+1)(d-2)}{d(d-1)(2d-1)}>0\).
\item
  \emph{Sub-solution.}
  \(\Pi(m,\allowbreak{}d):=\dfrac{2(d+m)}{2d-2+m}(g_m+1)-g_{m+1}=\dfrac{m(m-1)(2d^2-dm-d-2)}{(d-1)(2d+1-m)(2d+2-m)}\),
  which is \(\ge0\) for integers \(3\le m\le2d-2\) and \(<0\) for
  \(m\in\{2d-1,\allowbreak{}2d\}\).
\item
  \emph{Boundary indices.}
  \(q_{2d}\ge\prod_{k=d-2}^{2d-3}\frac{4d-3-k}{k+1}\) and
  \(q_{2d+1}\ge\prod_{k=d-2}^{2d-2}\frac{4d-2-k}{k+1}\); every factor is
  \(\ge1\) and the \(\lfloor d/2\rfloor+1\) factors with \(k\le3d/2-2\)
  are \(\ge5/3\). Hence
  \(q_{2d},\allowbreak{}q_{2d+1}\ge(5/3)^{\lfloor d/2\rfloor+1}\). Also
  \(g_{2d}=\frac{4d^3}{d-1}<g_{2d+1}=\frac{d(2d+1)(4d+1)}{d-1}\le9d^2\)
  for \(d\ge16\), and \((5/3)^{\lfloor d/2\rfloor+1}\ge9d^2\) for all
  \(d\ge36\) (false for \(d=34,\allowbreak{}35\)).
\end{enumerate}

\emph{Proof.} (i) Pascal's rule with the symmetry
\(\binom{T-1}{d-1+m}=\binom{T-1}{d-2}\) gives \(W_m=2W_{m-1}+2u_m\), and
\(u_{m+1}/u_m=(2d-2+m)/(d+m)\); divide. The value of \(q_3\) uses
\(W_2\) and \(u_3\) from Proposition 3.4. (ii) Direct: the numerator of
\(q_3-g_3\) over \(d(d-1)(2d-1)\) is
\((3d+1)(d+2)(2d-1)-3d^2(2d+3)=2d^2-3d-2\). (iii) The factorization is a
rational identity ({\small\texttt{induction\_\allowbreak{}check.\allowbreak{}py}}); for \(3\le m\le2d\)
the denominator is positive and the sign is that of \(2d^2-dm-d-2\),
which equals \(d-2\ge0\) at \(m=2d-2\) and \(-2\) at \(m=2d-1\), and is
decreasing in \(m\). (iv) For \(m=2d\), \(T-1=4d-3\) and the window
\(W_{2d-1}=\sum_{k=d-1}^{3d-2}\binom{4d-3}{k}\) contains the central
term \(k=2d-2\); keeping only it,
\(q_{2d}\ge\binom{4d-3}{2d-2}/\binom{4d-3}{d-2}=\prod_{k=d-2}^{2d-3}\frac{4d-3-k}{k+1}\)
by telescoping \(\binom{N}{j+1}/\binom{N}{j}=\frac{N-j}{j+1}\). A factor
is \(\ge1\) iff \(k\le2d-2\), and for \(k\le3d/2-2\) it is
\(\ge\frac{5d-2}{3d-2}\ge\frac53\); the number of integers
\(k\in[d-2,\allowbreak{}3d/2-2]\) is \(\lfloor d/2\rfloor+1\). The case \(m=2d+1\)
(\(T-1=4d-2\), central term \(k=2d-1\), factors
\(\ge\frac{5d}{3d-2}\ge\frac53\) for \(k\le3d/2-2\)) is identical.
\(g_{2d+1}\le9d^2\iff d^2-15d-1\ge0\iff d\ge16\). For the last claim,
\((5/3)^{19}>16400>9\cdot37^2=12321\) covers \(d=36,\allowbreak{}37\), and passing
from \(d\) to \(d+2\) multiplies the left side by \(5/3\) and the right
side by \(((d+2)/d)^2\le(38/36)^2<1.12\). \(\square\)

\emph{Proof of Theorem 6.1.} Since \(\Pi\ge0\) for \(m\le2d-2\) and
\(q\mapsto\frac{2(d+m)}{2d-2+m}(q+1)\) is increasing, (i)--(iii) give by
induction \(q_m>g_m\), i.e.~\(\hat c_{d,\allowbreak{}m}<0\), for \(3\le m\le2d-1\)
and every \(d\ge3\). For \(m\ge2d+2\), \(\hat c_{d,\allowbreak{}m}>0\) as noted. The
remaining indices are \(m\in\{2d,\allowbreak{}2d+1\}\) (i.e.~\(n\in\{2d-1,\allowbreak{}2d\}\)).
For \(d\ge36\), Lemma 6.2(iv) gives
\(\min(q_{2d},\allowbreak{}q_{2d+1})\ge(5/3)^{\lfloor d/2\rfloor+1}\ge9d^2\ge g_{2d+1}>g_{2d}\),
so \(\hat c_{d,\allowbreak{}2d}<0\) and \(\hat c_{d,\allowbreak{}2d+1}<0\), i.e.~\(n_d=2d\). For
\(3\le d\le35\), {\small\texttt{tail\_\allowbreak{}cases.\allowbreak{}py}} compares \(q_m\) and \(g_m\) in
exact integer arithmetic: \(q_{2d}>g_{2d}\) for all \(3\le d\le60\), and
\(q_{2d+1}>g_{2d+1}\) exactly for \(5\le d\le60\), failing at \(d=3,\allowbreak{}4\);
this gives \(n_d=2d\) for \(5\le d\le35\) and \(n_d=2d-1\) for
\(d=3,\allowbreak{}4\). Together with \(\hat c_{d,\allowbreak{}0}=\hat c_{d,\allowbreak{}1}=\hat c_{d,\allowbreak{}2}=0\)
(Proposition 3.4) and \(c_{d,\allowbreak{}n}\propto\hat c_{d,\allowbreak{}n+1}\), the sign pattern
is complete. For \(d=2\), (iii) is vacuous and
\(\hat c_{2,\allowbreak{}m}=(m-6)(2^{m+1}-2)+2m(m+4)\) is \(0\) at \(m=3\) and
positive for \(m\ge4\) (it is \(4,\allowbreak{}28,\allowbreak{}120,\allowbreak{}\dots\) and the first term is
\(\ge0\) from \(m=6\) on). \(\square\)

\emph{What the finite check covers.} Exactly the \(66\) integer
inequalities \(q_{2d}>g_{2d}\) and \(q_{2d+1}\gtrless g_{2d+1}\) for
\(3\le d\le35\); each is a comparison of two explicit rationals built
from binomial sums ({\small\texttt{fractions.\allowbreak{}Fraction}}, {\small\texttt{math.\allowbreak{}comb}}).
Independently, {\small\texttt{verify\_\allowbreak{}signs.\allowbreak{}py}} recomputes every \(c_{d,\allowbreak{}n}\)
from the three-term definition in exact rational arithmetic for
\(2\le d\le40\), \(n\le400\) and confirms the pattern and the values of
\(n_d\). The single central term already suffices analytically from
\(d\ge11\) (\(m=2d+1\)) and \(d\ge10\) (\(m=2d\)), so the check could be
shortened by summing a few more central terms; we did not pursue this.

\emph{Remark 6.3 (Descartes route).} Theorem 6.1 gives a second proof of
the one-zero statement of Theorem 4.1 (without the bounds (4.1)),
computer-assisted for \(3\le d\le35\). Write
\(\tilde N_d(\kappa)=\sum_{m\ge3}a_m\kappa^m\) (entire, by Lemma 2.1(i))
with \(a_m<0\) for \(3\le m\le n_d+1\) and \(a_m>0\) for \(m\ge n_d+2\)
(Theorem 6.1 with
\(\operatorname{sign}a_m=\operatorname{sign}c_{d,\allowbreak{}m-1}\)). Then
\[Q(\kappa):=\frac{\tilde N_d(\kappa)}{\kappa^{\,n_d+1}}=\sum_{3\le m\le n_d+1}a_m\kappa^{m-n_d-1}+\sum_{m\ge n_d+2}a_m\kappa^{m-n_d-1}\]
is a sum of negative multiples of nonpositive powers (the first sum,
which contains the negative constant \(a_{n_d+1}\)) and of positive
multiples of positive powers (the second sum). Termwise differentiation
(legitimate on \((0,\allowbreak{}\infty)\) for an entire function divided by a
monomial) gives \(Q'(\kappa)=\sum_m(m-n_d-1)a_m\kappa^{m-n_d-2}\), in
which every term is \(\ge0\) and the terms with \(m\ge n_d+2\) are
\(>0\): \(Q\) is strictly increasing on \((0,\allowbreak{}\infty)\). Since
\(n_d+1\ge6>3\), the term \(a_3\kappa^{2-n_d}\) forces
\(Q(\kappa)\to-\infty\) as \(\kappa\downarrow0\), and
\(Q(\kappa)\to+\infty\) as \(\kappa\to\infty\). Hence \(Q\), and so
\(\tilde N_d\) and \(H_d\), has exactly one positive zero, with
\(H_d<0\) before it and \(H_d>0\) after. (This is the standard argument
that a power series with one coefficient sign change has at most one
positive zero, and the closing step of 7H9F Theorem 8.2. Version 1 of
this note misdescribed \(Q\) as ``a negative constant plus a series with
positive coefficients'': the negative-power terms are present; they do
not spoil the monotonicity because their derivatives are positive.
{\small\texttt{h6\_\allowbreak{}remark63\_\allowbreak{}check.\allowbreak{}py}} checks \(Q'>0\) on a grid for
\(d=3,\allowbreak{}\dots,\allowbreak{}6\).)

\emph{Remark 6.4 (why the 7H9F window argument does not transfer).} In
7H9F the coefficient is a parity-truncated central window mean of the
form ``\(r-\mathbb E\,\allowbreak{}\delta^2\)'', controlled by a likelihood-ratio
ordering in the window size. Here Lemma 3.1 collapses the three-term
expression to a \emph{full} window mass \(W_{m-1}\) against a single
boundary binomial \(u_m\); the sign question becomes window mass versus
boundary mass, monotone enough for the sub-solution induction up to
\(m=2d-1\), with the last two indices exponentially safe.

\subsection{7. Numerical illustration}\label{numerical-illustration}

{\small\texttt{chain\_\allowbreak{}numerics.\allowbreak{}py}} (mpmath, 60 digits; 220-step bisection for
\(\kappa_f\) on \([\kappa_1,\allowbreak{}2d]\) with \(\kappa_1=2(d-2)(d+1)/d\), and
for \(\kappa_c\) on \((\kappa_f,\allowbreak{}20d+50]\)) gives the values below,
printed to 12--14 significant digits;
{\small\texttt{repro/\allowbreak{}chain\_\allowbreak{}values\_\allowbreak{}50digits.\allowbreak{}txt}} lists every quantity to 50
digits for fourteen values of \(d\). All rows satisfy the proved
relations \(\kappa_1\le\kappa_f<2d\), \(M_d(\kappa_f)\ge d(d-1)/2\),
\(b_d(\kappa_f)<1/2\), \(\kappa_f<\kappa_c\),
\(\lambda_{\min}<\lambda_c<\lambda_{0,\allowbreak{}d}\) (asserted in the script).

\begin{longtable}[]{@{}llll@{}}
\caption{the fold and the bounds (4.1).}\tabularnewline
\toprule\noalign{}
\(d\) & \(\kappa_1\) & \(\kappa_f\) & \(2d-\kappa_f\) \\
\midrule\noalign{}
\endfirsthead
\toprule\noalign{}
\(d\) & \(\kappa_1\) & \(\kappa_f\) & \(2d-\kappa_f\) \\
\midrule\noalign{}
\endhead
\bottomrule\noalign{}
\endlastfoot
3 & 2.66666666667 & 3.2327088362698 & 2.76729 \\
5 & 7.2 & 8.4315508759981 & 1.56845 \\
17 & 31.7647058824 & 33.907100295951 & 0.0928997 \\
100 & 197.96 & 199.99999999999 & \(9.68\cdot10^{-12}\) \\
\end{longtable}

\begin{longtable}[]{@{}
  >{\raggedright\arraybackslash}p{(\linewidth - 8\tabcolsep) * \real{0.2000}}
  >{\raggedright\arraybackslash}p{(\linewidth - 8\tabcolsep) * \real{0.2000}}
  >{\raggedright\arraybackslash}p{(\linewidth - 8\tabcolsep) * \real{0.2000}}
  >{\raggedright\arraybackslash}p{(\linewidth - 8\tabcolsep) * \real{0.2000}}
  >{\raggedright\arraybackslash}p{(\linewidth - 8\tabcolsep) * \real{0.2000}}@{}}
\caption{the fold radius and the minimal multiplier.}\tabularnewline
\toprule\noalign{}
\begin{minipage}[b]{\linewidth}\raggedright
\(d\)
\end{minipage} & \begin{minipage}[b]{\linewidth}\raggedright
\(M_d(\kappa_f)\)
\end{minipage} & \begin{minipage}[b]{\linewidth}\raggedright
\(d(d-1)/2\)
\end{minipage} & \begin{minipage}[b]{\linewidth}\raggedright
\(b_d(\kappa_f)\)
\end{minipage} & \begin{minipage}[b]{\linewidth}\raggedright
\(\lambda_{\min}\)
\end{minipage} \\
\midrule\noalign{}
\endfirsthead
\toprule\noalign{}
\begin{minipage}[b]{\linewidth}\raggedright
\(d\)
\end{minipage} & \begin{minipage}[b]{\linewidth}\raggedright
\(M_d(\kappa_f)\)
\end{minipage} & \begin{minipage}[b]{\linewidth}\raggedright
\(d(d-1)/2\)
\end{minipage} & \begin{minipage}[b]{\linewidth}\raggedright
\(b_d(\kappa_f)\)
\end{minipage} & \begin{minipage}[b]{\linewidth}\raggedright
\(\lambda_{\min}\)
\end{minipage} \\
\midrule\noalign{}
\endhead
\bottomrule\noalign{}
\endlastfoot
3 & 4.0410906 & 3 & 0.301630225423 & 5.3587282768762 \\
5 & 21.107032 & 10 & 0.43508469636 & 9.6895511914568 \\
17 & 3643.6645 & 136 & 0.498767685262 & 33.990875208891 \\
100 & \(1.0640\cdot10^{15}\) & 4950 & \(0.5-2.4\cdot10^{-14}\) &
\(200-1.9\cdot10^{-13}\) \\
\end{longtable}

\begin{longtable}[]{@{}lllll@{}}
\caption{the contact.}\tabularnewline
\toprule\noalign{}
\(d\) & \(\kappa_c\) & \(b_c\) & \(\lambda_c\) & \(\lambda_{0,\allowbreak{}d}\) \\
\midrule\noalign{}
\endfirsthead
\toprule\noalign{}
\(d\) & \(\kappa_c\) & \(b_c\) & \(\lambda_c\) & \(\lambda_{0,\allowbreak{}d}\) \\
\midrule\noalign{}
\endhead
\bottomrule\noalign{}
\endlastfoot
3 & 4.3442853031 & 0.400331536002 & 5.42585945949 & 6 \\
5 & 11.6440880518 & 0.573491039097 & 10.1519354776 & 15 \\
17 & 52.4714515166 & 0.676014299873 & 38.8094242433 & 153 \\
100 & 340.64385837 & 0.706438271107 & 241.099521574 & 5050 \\
\end{longtable}

\begin{longtable}[]{@{}llll@{}}
\caption{the contact point of the rate--distortion
function.}\tabularnewline
\toprule\noalign{}
\(d\) & \(D_c\) & \(R_c\) & \(\lambda_c/d\) \\
\midrule\noalign{}
\endfirsthead
\toprule\noalign{}
\(d\) & \(D_c\) & \(R_c\) & \(\lambda_c/d\) \\
\midrule\noalign{}
\endhead
\bottomrule\noalign{}
\endlastfoot
3 & 0.559823107522 & 0.579718136073 & 1.8086198 \\
5 & 0.53688642246 & 2.67111206245 & 2.0303871 \\
17 & 0.511063215405 & 16.6924477931 & 2.2829073 \\
100 & 0.495935519424 & 119.118709893 & 2.4109952 \\
\end{longtable}

The \(d=3\) and \(d=5\) rows agree to \(\le5\cdot10^{-9}\), and the
\(d=17\) row to all printed digits, with an independent brute-force
maximization of \(G_{d,\allowbreak{}\lambda}\) over a grid in \(b\) (no use of
\(H_d\), \(F_d\), \(\Psi_d\) or Lemma 3.1) performed by the adversarial
reviewer; the two-piece formula agrees with a brute-force evaluation of
the supremum in (A1) to \(10^{-14}\) or better at sample points on both
pieces. An external recomputation (AIRR workshop, 2026-09-08; 75-digit
arithmetic, roots of \(H_d\) and \(F_d\) evaluated through the Kummer
derivatives \(M_d'={}_1F_1(2;\allowbreak{}d+1;\allowbreak{}\kappa)/d\),
\(M_d''=2\,\allowbreak{}{}_1F_1(3;\allowbreak{}d+2;\allowbreak{}\kappa)/(d(d+1))\), i.e.~without Lemma 3.1 or
\(\Psi_d\)) reproduced \(\kappa_f\), \(\lambda_{\min}\), \(\kappa_c\),
\(b_c\), \(\lambda_c\), \(D_c\), \(R_c\) for \(d=3,\allowbreak{}5,\allowbreak{}17,\allowbreak{}100\) to its
printed 20 digits; {\small\texttt{workshop\_\allowbreak{}recheck.\allowbreak{}py}} repeats that route at
70 digits and agrees with the \(\Psi\)-route values of
{\small\texttt{chain\_\allowbreak{}values\_\allowbreak{}50digits.\allowbreak{}txt}} to 50 significant digits, and
confirms the located roots by direct quadrature of the tilted
\(\mathrm{Beta}(1,\allowbreak{}d-1)\) moments (residuals \(\le10^{-43}\) for
\(d\le17\) and \(\le10^{-11}\) for \(d=100\), where the quadrature
itself is the limiting error). At \(d=100\), \(\lambda_c/d=2.4110\)
versus \(\alpha_*-c_*\log d/d=2.4104\) (1D2Q Thm. 7.3,
\(\alpha_*=2.455407\ldots\), \(c_*=0.977136\ldots\)).

The v003 workshop also compared the Kummer-derivative values with direct
quadrature after dividing the integrand by its mode value, using
85-digit arithmetic and intervals split around the mode. For
\(d=3,\allowbreak{}5,\allowbreak{}17,\allowbreak{}100\), the differences in \(K_d\), \(K_d^{\prime}\) and
\(K_d^{\prime\prime}\), and both fold/contact residuals, were below
\(10^{-65}\). This scaling removes the tiny-integral accuracy loss of
the earlier \(d=100\) quadrature. These are numerical comparisons, not
certified interval error bounds. The exact inputs and outputs are
included in the revision evidence.

\subsection{8. Open questions}\label{open-questions}

\begin{enumerate}
\def\labelenumi{\arabic{enumi}.}
\tightlist
\item
  \emph{Higher rank.} For rank-\(r\) projectors on
  \(\mathrm{Gr}_{\mathbb C}(r,\allowbreak{}d)\) the analogue of Theorem A needs a
  global matrix-Bingham (HCIZ) spectral extremum at fixed Frobenius
  norm, known only for \(\mathrm{Gr}_{\mathbb C}(2,\allowbreak{}4)\) (61Y0) and in a
  weak/strong-field sense elsewhere; and the radial normalizer is a
  matrix-argument \({}_1F_1\) for which no first-order equation of the
  type of Lemma 3.1 is available. Whether the one-fold property persists
  is open.
\item
  \emph{Matrix-Jacobi normalizers.} Lemma 3.1 exists because \(M_d\) has
  a single numerator parameter; for \({}_1F_1(a;\allowbreak{}c;\allowbreak{}\kappa)\) with \(a>1\)
  (Jacobi-type overlap laws \(\mathrm{Beta}(a,\allowbreak{}c-a)\)) the reduction
  produces a second-order relation and the \((M,\allowbreak{}\kappa)\) transversality
  argument must be replaced by one in a three-dimensional phase space. A
  clean criterion in terms of \((a,\allowbreak{}c)\) for exactly one fold would unify
  7H9F (\(L_q\) = even part of \(M_{q+1}\)) and the present result.
\item
  \emph{General Bingham fixed-norm extremum.} The spectral input of
  Theorem A (one positive spike maximizes the projective Laplace
  transform at fixed traceless norm, all degrees) is the
  Brazitikos--Pandis extremum. A direct proof adapted to the tilted
  \(\mathrm{Beta}\) structure, or an extension to all rank-\(r\)
  Dirichlet partition functions, would remove the only external black
  box of this line.
\item
  Not addressed here, as in 1D2Q: a matching third-order coding theorem
  and the uniform double scaling \(d\to\infty\), \(D\downarrow0\).
\end{enumerate}

\subsection{9. Reproducibility}\label{reproducibility}

All scripts are in {\small\texttt{repro/\allowbreak{}}} next to this file, with a
{\small\texttt{README.\allowbreak{}md}} listing commands, expected final lines and wall times
(Python 3.12.6, sympy 1.14.0, mpmath 1.3.0; {\small\texttt{python\ run\_\allowbreak{}all.\allowbreak{}py}}
runs everything and writes {\small\texttt{runtimes.\allowbreak{}txt}}).

{\def\LTcaptype{none} % do not increment counter
\begin{longtable}[]{@{}
  >{\raggedright\arraybackslash}p{(\linewidth - 6\tabcolsep) * \real{0.2500}}
  >{\raggedright\arraybackslash}p{(\linewidth - 6\tabcolsep) * \real{0.2500}}
  >{\raggedright\arraybackslash}p{(\linewidth - 6\tabcolsep) * \real{0.2500}}
  >{\raggedright\arraybackslash}p{(\linewidth - 6\tabcolsep) * \real{0.2500}}@{}}
\toprule\noalign{}
\begin{minipage}[b]{\linewidth}\raggedright
script
\end{minipage} & \begin{minipage}[b]{\linewidth}\raggedright
certifies
\end{minipage} & \begin{minipage}[b]{\linewidth}\raggedright
arithmetic
\end{minipage} & \begin{minipage}[b]{\linewidth}\raggedright
time
\end{minipage} \\
\midrule\noalign{}
\endhead
\bottomrule\noalign{}
\endlastfoot
{\small\texttt{sym\_\allowbreak{}check.\allowbreak{}py}} & Lemma 3.2, identities (a)--(f) of Theorem 4.1,
\(\hat c_{d,\allowbreak{}3}\), with \(d\) symbolic & sympy exact & 2.0 s \\
{\small\texttt{check\_\allowbreak{}formula.\allowbreak{}py}} & \(N_d=\frac{d-1}{d\kappa}\tilde N_d\) and
(3.3) coefficientwise, \(2\le d\le12\), \(n\le60\) & {\small\texttt{Fraction}} &
0.3 s \\
{\small\texttt{verify\_\allowbreak{}signs.\allowbreak{}py}} & (3.3) against the three-term definition and
the sign pattern of Theorem 6.1, \(2\le d\le40\), \(n\le400\) &
{\small\texttt{Fraction}} & 80.6 s \\
{\small\texttt{induction\_\allowbreak{}check.\allowbreak{}py}} & factorization of \(\Pi(m,\allowbreak{}d)\) and its
sign table, \(3\le d\le60\) & sympy & 3.3 s \\
{\small\texttt{tail\_\allowbreak{}cases.\allowbreak{}py}} & recursion (i) for \(d\le29\); the boundary
comparisons of Theorem 6.1 for \(3\le d\le60\); the central-term and
\((5/3)^{\lfloor d/2\rfloor+1}\) bounds & exact integers & 0.1 s \\
{\small\texttt{boundary\_\allowbreak{}exact.\allowbreak{}py}} & the 66 boundary comparisons as explicit
rationals \(q_m\), \(g_m\), \(\hat c_{d,\allowbreak{}m}\) (written to
{\small\texttt{boundary\_\allowbreak{}exact.\allowbreak{}txt}}) & {\small\texttt{Fraction}} & 0.1 s \\
{\small\texttt{h6\_\allowbreak{}remark63\_\allowbreak{}check.\allowbreak{}py}} & expansion (2.5) with \(d,\allowbreak{}\lambda\)
symbolic; \(G_{2,\allowbreak{}3}(b)=-\frac9{20}b^4+\frac9{35}b^6+O(b^8)\); \(Q'>0\)
in Remark 6.3 for \(d=3,\allowbreak{}\dots,\allowbreak{}6\) & sympy; mpmath 30 digits & 2.9 s \\
{\small\texttt{chain\_\allowbreak{}numerics.\allowbreak{}py}} & Section 7 tables (14 values of \(d\)), the
asserted inequalities, {\small\texttt{chain\_\allowbreak{}values\_\allowbreak{}50digits.\allowbreak{}txt}} & mpmath, 60
digits & 1.4 s \\
{\small\texttt{workshop\_\allowbreak{}recheck.\allowbreak{}py}} & the 66 signs from the three-term
definition; Section 7 values by the Kummer-derivative route (70 digits)
against the workshop's 20 digits and
{\small\texttt{chain\_\allowbreak{}values\_\allowbreak{}50digits.\allowbreak{}txt}}; quadrature residuals &
{\small\texttt{Fraction}}; mpmath & 2.2 s \\
\end{longtable}
}

No proof step depends on floating point. The proof of Theorem 4.1 and
Theorem 5.1 uses only Lemmas 2.1, 3.1--3.2, the rational identities
(a)--(f) (which are one-line hand computations, machine-checked for
safety), and Theorem A. Theorem 6.1 depends on {\small\texttt{tail\_\allowbreak{}cases.\allowbreak{}py}}
for \(3\le d\le35\) as stated; the 66 exact rationals it compares are
listed in {\small\texttt{boundary\_\allowbreak{}exact.\allowbreak{}txt}}, so the finite part of that proof
can be checked by hand or by any exact-arithmetic system.

\subsection{Changes in internal revision v004 (9 September
2026)}\label{changes-in-internal-revision-v004-9-september-2026}

\begin{enumerate}
\def\labelenumi{\arabic{enumi}.}
\tightlist
\item
  Documented and repaired a local arithmetic error in the proof of the
  external power-sum inequality {[}BP{]}, with the exact arXiv version
  and page identified. No main theorem statement changed.
\item
  Corrected the reproduction runner to propagate subprocess failures and
  made the expected symbolic identities and coefficient signs mandatory
  assertions. The nine programs were executed successfully in the audit
  supplement; the 66 proof-critical signs were also reconstructed
  separately with exact arithmetic.
\item
  Preserved the earlier deposited v003 and all its evidence. The current
  PDF is a new revision prepared for formal assessment, not an editorial
  acceptance.
\end{enumerate}

\subsection{Changes in internal revision v003 (8 September
2026)}\label{changes-in-internal-revision-v003-8-september-2026}

\begin{enumerate}
\def\labelenumi{\arabic{enumi}.}
\tightlist
\item
  Added a proof interface for the imported envelope: the precise
  centered power-sum hypothesis, the Haar/Dirichlet Laplace comparison,
  the entropy lower bound and active-channel attainment. The cited
  Brazitikos-Pandis extremal inequality remains an external theorem.
\item
  Corrected (A3): forgetting an independent mixture flag gives an
  information upper bound; the matching converse supplies equality on
  the frontier.
\item
  Replaced an imprecise moving-target limit notation by the
  corresponding difference.
\item
  Kept table captions with their tables and added a scaled-quadrature
  check that avoids tiny-integral loss at \(d=100\); results and limits
  are recorded separately.
\item
  Qualified historical AI-review and computation provenance.
\end{enumerate}

\subsection{Changes in version 2 (8 September
2026)}\label{changes-in-version-2-8-september-2026}

Version 2 answers the evaluation of the AIRR workshop of 8 September
2026 (its items for this note are numbered W1--W4 below). Version 1 is
kept as {\small\texttt{paper\_\allowbreak{}v1.\allowbreak{}md}}.

\begin{enumerate}
\def\labelenumi{\arabic{enumi}.}
\tightlist
\item
  (W1) \emph{Remark 6.3.} Version 1 claimed that
  \(\tilde N_d(\kappa)/\kappa^{n_d+2}\) is ``a negative constant plus a
  series with positive coefficients''. This is false: after division
  there are negative coefficients on negative powers, a positive
  constant and positive coefficients on positive powers. The remark now
  divides by \(\kappa^{n_d+1}\), differentiates termwise (every term of
  \(Q'\) is \(\ge0\), the negative-power terms contributing positive
  derivatives), and uses the left limit \(-\infty\); the conclusion (one
  zero, one sign change) is unchanged, and Theorem 4.1 never depended on
  it. Checked on a grid by {\small\texttt{h6\_\allowbreak{}remark63\_\allowbreak{}check.\allowbreak{}py}}.
\item
  (W2) \emph{Fact (H6).} Version 1 stated ``\(0\) is a strict local
  maximum iff \(\lambda<\lambda_{0,\allowbreak{}d}\)'' for all \(d\ge2\). At \(d=2\),
  \(\lambda=\lambda_{0,\allowbreak{}2}=3\) this is false: \(\mu_{3,\allowbreak{}2}=0\) and
  \(G_{2,\allowbreak{}3}(b)=-\frac9{20}b^4+\frac9{35}b^6+O(b^8)\), so \(0\) remains a
  strict local maximum, as Corollary 5.2 already implied. (H6) now
  separates \(\lambda<\lambda_{0,\allowbreak{}d}\), \(\lambda>\lambda_{0,\allowbreak{}d}\) and
  \(\lambda=\lambda_{0,\allowbreak{}d}\), the last split into \(d\ge3\) (cubic term,
  not a local maximum) and \(d=2\) (quartic term, strict local maximum),
  and records which cases are used. The expansion is now derived in
  Lemma 2.1 (new) and checked symbolically.
\item
  (W3) \emph{Table and title page.} The single twelve-column table of
  Section 7 overlapped in the PDF; it is replaced by four tables of at
  most five columns (Tables 1--4), with the values recomputed at 60
  digits and listed to 50 digits in
  {\small\texttt{repro/\allowbreak{}chain\_\allowbreak{}values\_\allowbreak{}50digits.\allowbreak{}txt}}. The bisection of
  {\small\texttt{chain\_\allowbreak{}numerics.\allowbreak{}py}} was replaced by a fixed 220-step bisection
  because
  \texttt{mpmath.findroot(...,\ \textquotesingle{}bisect\textquotesingle{})}
  stopped at about 28 correct digits, short of the 50 digits announced
  in version 1 (the printed 10-digit values were unaffected). The
  duplicated date on the title page is removed (the date is set by the
  rendering metadata only).
\item
  (W4) \emph{Framing and self-containment.} The introduction now says
  explicitly that the note is standalone, which single statement it
  imports (Theorem A), that Lemma 3.1 is a contiguous relation of
  Kummer's function whose novelty lies in its use, and what a reader
  gains without 1D2Q. Lemma 2.1 (new) reproves the facts previously
  cited as 1D2Q eqs. (10), (11), (15); the direct reason for
  \(R_d(0)=+\infty\) is added in the proof of Theorem 5.1(f). The note
  is deliberately not merged with 1D2Q.
\item
  \emph{Verification of the workshop's checks.} Its 66 boundary sign
  comparisons and its 75-digit fold/contact values for \(d=3,\allowbreak{}5,\allowbreak{}17,\allowbreak{}100\)
  were recomputed by routes independent of both the workshop's script
  and version 1's scripts ({\small\texttt{workshop\_\allowbreak{}recheck.\allowbreak{}py}},
  {\small\texttt{boundary\_\allowbreak{}exact.\allowbreak{}py}}); all agree (Section 7).
\end{enumerate}

No theorem statement changed. The proofs of Theorems 4.1, 5.1 and 6.1
are unchanged except for the explicit case split in (H6), which they use
only in the cases \(\lambda<\lambda_{0,\allowbreak{}d}\), \(\lambda>\lambda_{0,\allowbreak{}d}\)
and (\(d\ge3\), \(\lambda=\lambda_{0,\allowbreak{}d}\)).

\subsection{AI-assistance statement}\label{ai-assistance-statement}

The received version declares substantial generative-AI assistance in
drafting, mathematical derivations, code and review, attributed there to
``Claude Fable 5.1'' (Anthropic). That historical model identity, the
reported review sessions and their independence have not been
authenticated by this workshop. Historical descriptions of an ``author''
or ``reviewer'' script identify the supplied files, not a verified
independent editorial review. OpenAI Codex assisted with the present
private workshop revision: checking arguments and sources, correcting
identified errors, running the checks recorded in the accompanying
revision report, and preparing the PDF. Numerical checks support only
their tested instances. This assistance is not an independent human peer
review or a formal AIRR model-assessment report. Lluis Eriksson is the
declared author and AIRR founder/operator; this conflict must be handled
by the separate editorial process. Scientific authorship and approval
remain with the author.

\subsection{References}\label{references}

{[}1D2Q{]} L. Eriksson, \emph{Exact Rate--Distortion Theory for
Complex-Projective Born Prediction: Finite-Dimensional Bingham
Frontiers, Thermodynamic Coexistence, and Worst-State Capacity},
ARR-2026-1D2QYXPCVY9H7ANB, v1, 2026-08-13.

{[}61Y0{]} L. Eriksson, \emph{Complete Rank-Two Born-Prediction
Rate--Distortion on \(\mathrm{Gr}_{\mathbb C}(2,\allowbreak{}4)\): All-Field
Matrix--Bingham Rigidity and a Unique Coexistence Transition},
ARR-2026-61Y0FFA39M8KMBJ5, v1, 2026-08-13.

{[}7H9F{]} L. Eriksson, \emph{Complete Rate--Distortion Phase Diagram of
Haar Oriented Two-Planes: One-Change Hypergeometric Coefficients, Unique
Coexistence, and No Reentrance}, ARR-2026-7H9FAPTBZA897AMJ, v2,
2026-08-14.

{[}BP{]} S. Brazitikos and C. Pandis, \emph{Sharp Inequalities for
Symmetric Polynomials, Hunter's Conjecture, and Moments of Exponential
Random Variables}, arXiv:2512.12254v1 {[}math.PR{]}, 2025 (Proposition
5.5; local estimate clarified in Section 2).

{[}DLMF{]} NIST Digital Library of Mathematical Functions, Chapter 13
(Confluent Hypergeometric Functions), §13.3 (recurrence relations and
derivatives), https://dlmf.nist.gov/13.3.

{[}DC{]} A. Dytso and M. Cardone, \emph{Uniform Distribution on
\((n-1)\)-Sphere: Rate-Distortion under Squared Error Distortion},
arXiv:2401.04248 {[}cs.IT{]}, 2024.

{[}EE{]} T. Eisele and R. S. Ellis, \emph{Multiple phase transitions in
the generalized Curie--Weiss model}, J. Stat. Phys. 52 (1988), no. 1--2,
161--202.

{[}KS{]} D. Karp and S. M. Sitnik, \emph{Log-convexity and log-concavity
of hypergeometric-like functions}, J. Math. Anal. Appl. 364 (2010),
384--394.

{[}KK{]} S. I. Kalmykov and D. B. Karp, \emph{Log-convexity and
log-concavity for series in gamma ratios and applications},
arXiv:1211.2882 {[}math.CA{]}, 2012.

{[}S{]} S. M. Sitnik, \emph{A conjecture on monotonicity of a ratio of
Kummer hypergeometric functions}, arXiv:1207.0936 {[}math.CA{]}, 2012.

\end{document}
